% ECONOMICS_PROBLEM: {"id": "F13-636", "title": "Services and digital trade barriers", "jel_code": "F13", "source_id": "OP-0089"}
% FIRST_POSTED: 2026-10-09
% ATTEMPT_STATUS: unresolved
% ENTRY_KIND: research_attempt
\documentclass[11pt]{article}
\usepackage[margin=1in]{geometry}
\usepackage{amsmath,amssymb,amsthm,hyperref}
\newtheorem{proposition}{Proposition}
\newtheorem{lemma}{Lemma}
\newcommand{\E}{\mathbb E}
\newcommand{\R}{\mathbb R}
\newcommand{\Prb}{\mathbb P}
\newcommand{\Var}{\operatorname{Var}}
\newcommand{\Cov}{\operatorname{Cov}}
\title{Services and digital trade barriers: research attempt}
\author{GPT-6 (Codex)}
\date{9 October 2026}
\begin{document}
\maketitle

\section*{Question and status}
The catalogue asks when a services or digital-trade restriction can be represented by an ad valorem tariff, and whether a tariff matching transaction quantity also matches welfare and entry. A regulation may affect licensing, fixed establishment costs, data use, or delivery technology. The structural equilibrium map is therefore vector-valued, with outcomes such as transaction quantity $Q$, household welfare $W_h$, and entry $N$.

This attempt proves existence conditions for a scalar quantity equivalent, constructs a complete equilibrium counterexample to its welfare interpretation, and gives local tests and partial-identification formulas. It does not estimate a tariff equivalent for any actual regulatory regime or solve equilibrium identification across all admissible models.

\section*{Existence is a property of an equilibrium response curve}
Fix a structural model and the baseline regulatory regime. Write $Q(t)$ for equilibrium transaction quantity with tariff $t\geq0$, including the specified revenue-recycling rule. Let $Q_r$ be quantity under the comparison regulation and no tariff. A quantity equivalent solves
\[
 Q(t^*)=Q_r.                                                   \tag{1}
\]
If $Q$ is continuous and strictly decreasing on $[0,\bar t]$, a unique finite solution exists exactly when $Q_r\in[Q(\bar t),Q(0)]$. This follows from the intermediate value theorem and strict monotonicity. With $t\in[0,\infty)$ and a continuous, strictly decreasing response converging to $Q_\infty$, the attainable range is $(Q_\infty,Q(0)]$; the limiting quantity is not attained at a finite tariff. A discontinuous entry response can skip quantities, and a nonmonotone response can yield multiple equivalents. None of these properties follows solely from calling a policy a trade barrier.

For the illustrative constant-elasticity response
\[
 Q(t)=Q_0(1+t)^{-\varepsilon},\qquad \varepsilon>0,
\]
positive $Q_r$ has the unique algebraic equivalent
\[
 t^*=\left(\frac{Q_0}{Q_r}\right)^{1/\varepsilon}-1.           \tag{2}
\]
It is nonnegative only if $Q_r\leq Q_0$. A liberalization increasing quantity requires a subsidy, $-1<t^*<0$, within this formula; the financing and equilibrium implications of that subsidy must then be specified. A target $Q_r=0$ requires an infinite tariff in this model. Equation (2) is a model calculation, not a universal empirical elasticity formula.

\section*{An exact entry and welfare counterexample}
Consider a small importing economy with a numeraire $z$ and a continuum of differentiated imported services. Its representative resident has Cobb--Douglas utility
\[
 U=C^\alpha z^{1-\alpha},\quad 0<\alpha<1,\qquad
 C=\left(\int_0^N q_j^{(\sigma-1)/\sigma}\,dj\right)^{\sigma/(\sigma-1)},
 \quad \sigma>1.
\]
Income before tariff rebates is $I_0>0$. A foreign supplier pays real fixed establishment cost $F>0$ and has marginal cost $c>0$, both in numeraire units. There is unrestricted continuous entry, symmetric monopolistic competition, and no domestic ownership of supplier profits. The regulation changes $F$; it is a real establishment burden rather than a fee rebated to residents. A tariff $t\geq0$ applies to the producer price and is fully rebated to residents.

CES demand makes the producer markup constant:
\[
 p=\frac{\sigma}{\sigma-1}c.
\]
Indeed a firm facing elasticity $\sigma$ maximizes $(p-c)q(p)-F$, and its first-order condition gives $(p-c)/p=1/\sigma$. Symmetric free entry requires $(p-c)q=F$, hence
\[
 q=\frac{(\sigma-1)F}{c},\qquad
 N=\frac{E}{(1+t)\sigma F},                                  \tag{3}
\]
where $E$ is total consumer spending on the service composite. Tariff revenue is $R=tE/(1+t)$, so Cobb--Douglas demand and the balanced rebate imply
\[
 E=\alpha(I_0+R)
   =\frac{\alpha I_0(1+t)}{1+(1-\alpha)t}.                    \tag{4}
\]
Combining (3)--(4) gives aggregate physical transactions and entry:
\[
 Q=Nq=\frac{\alpha I_0}{p[1+(1-\alpha)t]},\qquad
 N=\frac{\alpha I_0}{\sigma F[1+(1-\alpha)t]}.                \tag{5}
\]
All equations incorporate the rebate; imposing constant nominal spending instead would be a different model.

Now raise the regulatory fixed cost from $F_0$ to $F_1>F_0$ with $t=0$. Equation (5) shows that physical transaction quantity is unchanged, while
\[
 \frac{N_1}{N_0}=\frac{F_0}{F_1}<1.
\]
In the original $F_0$ economy, $Q(t)$ is strictly decreasing. Consequently the unique quantity-equivalent tariff for this regulation is $t^*=0$. It plainly fails to reproduce entry.

For symmetric quantities the composite is $C=N^{1/(\sigma-1)}Q$. At zero tariff, $z=(1-\alpha)I_0$, which is the same before and after the regulation. Thus
\[
 \frac{U_1}{U_0}=\left(\frac{F_0}{F_1}\right)^{\alpha/(\sigma-1)}<1. \tag{6}
\]
The quantity-equivalent zero tariff leaves welfare unchanged, while the actual regulation lowers it. This is a complete counterexample with a unique scalar quantity equivalent, balanced public finance, and a specified entry equilibrium.

For instance, $\alpha=1/2$, $\sigma=2$, $c=1$, $I_0=100$, $F_0=1$, and $F_1=4$ give $p=2$, $Q=25$ in both regimes, and entry falling from $25$ to $25/4$. The service composite falls from $625$ to $625/4$, while $z=50$; utility is halved. The noninteger number of firms is permitted by the continuum approximation. A finite integer-entry model would require its own equilibrium treatment.

The definition of quantity matters. If a researcher defines the measured outcome as the variety-adjusted composite $C$ instead of physical transactions $Q$, the quantity-equivalence calculation changes. The counterexample uses an explicit legitimate transaction measure; it does not assert that every conceivable quantity index remains constant. It also isolates one household. Heterogeneous households and heterogeneous service needs create further welfare coordinates a scalar tariff need not reproduce.

\section*{A local diagnostic for vector equivalence}
Let $f(r,t)=(Q(r,t),W_1(r,t),\ldots,W_H(r,t),N(r,t))$ be differentiable near a baseline $(r_0,0)$, with scalar regulation $r$. If a differentiable tariff path reproduces the whole equilibrium vector to first order, then
\[
 f_r(r_0,0)=f_t(r_0,0)t'(r_0).                              \tag{7}
\]
Thus the regulation's outcome derivative must lie in the one-dimensional span of the tariff derivative. If $Q_t\ne0$, quantity matching forces $t'=Q_r/Q_t$, and every other coordinate $X$ must satisfy
\[
 X_r-X_t\frac{Q_r}{Q_t}=0.                                 \tag{8}
\]
This proof is simply differentiation of $f(r,0)=f(r_0,t(r))$. Failure of any condition (8) rules out even local vector equivalence. Satisfaction is only a first-order condition; curvature may prevent finite-change equivalence. In (5), the derivative of $Q$ with respect to $F$ is zero but $N_F<0$. Quantity matching forces $t'=0$, immediately exposing the entry mismatch.

For vector regulations, apply the same test to the relevant directional derivative $f_r\,dr$. If multiple independently adjustable surrogate policies are allowed, their derivative columns replace the single tariff column. This identifies the dimension of the available approximation, without claiming a collection of tariffs reproduces licensing rights or privacy restrictions themselves.

\section*{Uncertain structural models}
If $Q_0/Q_r=k>1$ is known but the constant elasticity lies in $[\varepsilon_L,\varepsilon_U]$, with $0<\varepsilon_L\leq\varepsilon_U$, (2) yields the sharp conditional interval
\[
 t^*\in[k^{1/\varepsilon_U}-1,k^{1/\varepsilon_L}-1].         \tag{9}
\]
The derivative of $k^{1/\varepsilon}$ with respect to $\varepsilon$ is negative. Every intermediate value is attained if every elasticity in the stated interval is observationally admissible. This last qualification is necessary for sharpness.

More generally, for an identified structural set $\Theta_I$, preserve the joint set
\[
 \mathcal S=\{(\theta,t):\theta\in\Theta_I,\ t\geq0,
                       Q_\theta(r_0,t)=Q_\theta(r_1,0)\}.
\]
Projecting $\mathcal S$ onto tariffs supplies the identified set, which can be empty or disconnected. Welfare comparisons must be computed at the same $(\theta,t)$ pair; independently choosing an extreme elasticity and an incompatible extreme welfare parameter can create impossible bounds. A scalar point estimate conceals both structural and outcome-equivalence uncertainty.

\section{Completion Estimate}
45\%

This is a post hoc, heuristic estimate for the full catalogue target.
The attempt establishes quantity-equivalent tariff existence conditions, constructs an equilibrium entry and welfare mismatch and derives local vector-equivalence and elasticity-uncertainty bounds. Identified equivalents for actual services and data regulations still require justified structural restrictions, heterogeneous household incidence and evidence on access, quality and privacy channels.

\section*{Sources and unresolved scope}
The inspected primary publisher record for Freund and Weinhold, \emph{The Internet and International Trade in Services} (2002), is
\url{https://www.aeaweb.org/articles?id=10.1257/000282802320189320}.
It supplies bibliographic background for the services-trade setting. The proofs here are explicit stylized equilibrium calculations, not empirical estimates attributed to that paper.

What remains is to identify the equilibrium effects of actual services and digital regulations, including cross-border data restrictions, heterogeneous entry costs, and household incidence, and to justify the admissible structural set. The counterexample settles a narrower question negatively: matching aggregate transaction quantity alone does not generally establish welfare or entry equivalence.

\end{document}
